\documentclass[11pt,a4paper]{ltjsarticle}
\usepackage{amsmath,amssymb}
\usepackage[margin=24mm]{geometry}

\title{平成16年 B問題}
\author{数学科筋トレ部}
\date{}

\begin{document}
\maketitle

\section*{B 第1問}
有限次元線形空間 \(A,B,C,D\) の次元をそれぞれ \(a,b,c,d\) とし，線形写像
\(g:B \to C\) の階数を \(r\) とする。線形写像
\[
  F:\operatorname{Hom}(A,B)\otimes \operatorname{Hom}(C,D)
    \longrightarrow \operatorname{Hom}(A,D)
\]
を
\[
  F(f\otimes h)=hgf
\]
で定めたとき，\(F\) の階数を求めよ。

\section*{B 第2問}
\(n\) を自然数とし，\(L=\mathbb{C}(T_1,\ldots,T_n)\) を \(n\) 変数有理関数体とする。
\(\mathbb{C}\) 上の \(L\) の自己同型 \(\sigma\) を
\[
  \sigma(T_i)=T_{i+1}\quad i=1,\ldots,n-1,
  \qquad
  \sigma(T_n)=T_1
\]
で定める。\(K=\{f\in L\mid \sigma(f)=f\}\) を不変部分体とする。

\begin{enumerate}
  \item 拡大次数 \([L:K]\) を求めよ。
  \item \(L\) の線形部分空間 \(\bigoplus_{i=1}^{n}\mathbb{C}T_i\) を \(\sigma\) の作用に関する固有空間に分解せよ。
  \item 多項式 \(f_1,\ldots,f_n \in \mathbb{C}[T_1,\ldots,T_n]\) で \(K=\mathbb{C}(f_1,\ldots,f_n)\) となるようなものを一組与えよ。
  \item \(n=6\) のとき \(L\) と \(K\) の中間体をすべて求めよ。
\end{enumerate}

\section*{B 第3問}
複素 \(2\times2\) 行列の全体を \(M_2(\mathbb{C})\) で表す。
\[
  X=\{(A,B)\in M_2(\mathbb{C})\times M_2(\mathbb{C})
    \mid A^2=B^4=I,\ AB=B^3A\}
\]
とおき，\(H=\operatorname{GL}(2,\mathbb{C})\) の \(X\) への作用を
\[
  h(A,B)=(hAh^{-1},hBh^{-1}),
  \qquad h\in H,\ (A,B)\in X
\]
で定義する。

\begin{enumerate}
  \item \(X\) の中の \(H\) 軌道の個数を求めよ。
  \item \(X\) の中の各 \(H\) 軌道は
  \(M_2(\mathbb{C})\times M_2(\mathbb{C})\simeq \mathbb{C}^8\) の閉部分集合であることを示せ。
  ただし \(\mathbb{C}^8\) には \(16\) 次元ユークリッド空間としての位相を入れておく。
  \item 上記の各 \(H\) 軌道に対し，その一点 \((A,B)\) の固定群
  \[
    \{h\in H\mid h(A,B)=(A,B)\}
  \]
  の複素多様体としての次元を求めよ。
\end{enumerate}

\section*{B 第4問}
\(p\) を \(3\) 以上の素数とする。単位元をもつ可換環 \(A\) の可逆元のなす群 \(A^\times\) は，
位数 \(p^2\) の巡回群にはならないことを証明せよ。

\end{document}
